\section{Comparison and mathematical scope}\label{sec:comparison}
\subsection{Neighboring theories}
Static quadratic quantization supplies the optimal finite-center problem and its mass/cover bounds~\cite{graf}. Conditional projection is standard probability, and the convex perspective used in Section~\ref{sec:defect} is a standard convexity-preserving operation~\cite{boyd}. The broader relation between convex information, scoring rules and Bregman defects is classical~\cite{gneiting}. None of these ingredients, by itself, identifies the cost of reusing an outgoing label across distinct incoming posterior rows.

Borkar, Mitter and Tatikonda~\cite{borkar} formulate sequential vector quantization of a stationary Markov source as a partially observed stochastic control problem and characterize optimal encoding schemes by dynamic programming. Foster, Gray and Dunham~\cite{foster} use a finite-state quantizer whose old state selects a codebook and whose current codeword helps determine the next state. These are direct predecessors of the finite-state viewpoint. The present observer receives a noisy instrument report rather than the current source or exact posterior. Theorem~\ref{thm:main} gives a lower against every legal hard-state machine and a matching implementation with its internal deterministic deadline counted; it is not a claim that finite-state quantization or its control formulation is new.

Recursive marginal quantization~\cite{pages} studies accumulated discretization error in recursively quantized Markov dynamics, including the Euler scheme of a diffusion. Its one-step error propagation is close in spirit to~\eqref{eq:pathwise-recursion}. Here the exact normalized defect identifies the missing old information, and the separate pooling term quantifies compatibility across predecessor labels. The acquired-mass converse and complete clock charge are additional assertions; the recursive upper alone is not presented as a new general quantization principle.

Nonanticipative and sequential rate-distortion theories~\cite{stavrou2014,stavrou2018,tanaka} treat causal reproduction laws, information-rate lower bounds and, in Gaussian settings, filtering realizations and finite-horizon allocation. Their charged object is information rate, not a fixed total alphabet containing control phases and halt labels. In particular an information-rate allocation is not substituted for Lemma~\ref{lem:clock}. Conversely the present finite-state theorem does not provide a coding theorem or the semidefinite characterization of Gaussian sequential rate-distortion.

Atar and Zeitouni~\cite{atar} use projective contraction in nonlinear-filter stability. The metric mechanism in the raw realizations is part of that mature tradition; its elementary estimates are included to fix the constants and observation timing. Kara and Y\"uksel~\cite{kara} establish near-optimal finite-memory feedback policies through filter-stability methods in partially observed control. Ghomi, Linder and Y\"uksel~\cite{ghomi} obtain stationary optimal zero-delay codes and finite-window approximation for linear vector Markov sources under an infinite-horizon average criterion. The present results instead fix the physical report protocol, charge an exact internal deadline, and derive the hard-label exponent from the actual conditional report image. They do not replace those control or infinite-horizon conclusions.

\begin{center}\small
\begin{tabular}{p{.20\textwidth}p{.24\textwidth}p{.22\textwidth}p{.23\textwidth}}
\toprule
Theory & Observed object & Charged memory/information & Main kind of conclusion\\
\midrule
Sequential vector coding~\cite{borkar} & Markov source at encoder & Causal code/control structure & Optimal encoding characterization\\[3pt]
Finite-state vector coding~\cite{foster} & Current source vector and code state & State-selected codebooks & Design and storage/performance analysis\\[3pt]
Recursive quantization~\cite{pages} & Recursively approximated dynamics & Marginal codebooks & Accumulated approximation error\\[3pt]
Nonanticipative RDF~\cite{stavrou2018,tanaka} & Causal source/reproduction laws & Information rate & Distortion bounds and filtering design\\[3pt]
This theorem & Noisy positive-instrument report and old label & All persistent phases, labels and halt states & Exact defect, all-observer converse and matching order\\
\bottomrule
\end{tabular}
\end{center}
The table compares information patterns and assertions, not journal rank or historical priority. The published versions are cited where verified; absence of an exact bibliographic match is not a priority certificate.

\subsection{Relation to the retained foundations results}
The complete earlier articles are separate companions, not prerequisites to reading the present proof. The earlier posterior-pooling theorem already made conditional barycenters and Jensen losses operational for controlled noisy experiments. The earlier continuation-curvature theorem already obtained a charged-clock $(M-n)^{-2/d}$ law under a Wasserstein continuation certificate, including a stronger history-dependent action comparison than is claimed here. The Gaussian log-odds density estimate, positive congruence inverse, elementary erased-bit/offset devices, and integer phase allocation are inherited tools and have been explicitly re-proved or credited where used.

The present extension does not claim a new exponent merely because the same expression occurs again. It supplies the likelihood-normalized one-step perspective and its exact row/pooling converse, then obtains matching laws through a different same-report posterior-transport certificate. That certificate covers every positive invertible finite-state transition and every positive amount of the stated quantum improvement, including parameters outside the older sufficient $3\kappa$ block criterion. The rank-changing mixture combines that identity with a weighted acquired-stratum profile. The retained-bit theorem determines the additional autonomous phase bottleneck in the same scored experiment. The earlier relative-Haar theorem remains independently useful: it treats isometric updates for which summable stability fails and gives a different exact online law.

The General Theta problem is the connection between executable tests, causal implementation, acquired geometry and resources. Theorem~\ref{thm:main} proves that connection for a specified finite affine test class with informative likelihoods. It is neither a classification of all physical experiments nor a replacement of foundations by a single realization. The two raw models are independent verifications of the same general theorem; they are not premises in its proof.

Two substantial mathematical targets remain. First, one needs matching converses and implementable recursions for nonsummable, state-dependent continuation sensitivity, without inserting a global compatible cover as an assumption. Second, one needs a joint learning-and-memory law for unknown kernels under an explicitly fixed action and calibration interface. Unbounded stopping, arbitrary nondominated supports and fully optimal computation resources also require separate proofs. The finite-horizon results here do not acquire those conclusions by notation or by a limiting argument without additional uniform estimates.
